\IfFileExists{stacks-project.cls}{%
\documentclass{stacks-project}
}{%
\documentclass{amsart}
}
\usepackage{amssymb}

% For dealing with references we use the comment environment
\usepackage{verbatim}
\usepackage{fancyvrb}
\usepackage{pmboxdraw}
\newenvironment{reference}{\comment}{\endcomment}
%\newenvironment{reference}{}{}
\newenvironment{slogan}{\comment}{\endcomment}
\newenvironment{history}{\comment}{\endcomment}

% For source figures retained as assets
\usepackage{graphicx}

% For commutative diagrams we use Xy-pic
\usepackage[all]{xy}

% We use 2cell for 2-commutative diagrams.
\xyoption{2cell}
\UseAllTwocells

% We use multicol for the list of chapters between chapters
\usepackage{multicol}

% This is generally recommended for better output
\usepackage{lmodern}
\usepackage[T1]{fontenc}
\usepackage{textalpha}

% For cross-file-references
% We prefer the snapshot tmp/xref/#2.aux,
% which latexmkrc replaces atomically, so
% that documents compiled simultaneously
% never read a half-written .aux file.
% The last argument keeps links to #2.pdf.
\usepackage{xr-hyper}
\newcommand{\maybeexternaldocument}[2]{%
  \IfFileExists{#2.aux}{%
    \externaldocument[#1]{#2}%
  }{%
    \PackageWarning{stack}{Missing #2.aux; external references with prefix #1 unavailable}%
  }%
}
\let\externaldocumentorig\externaldocument
\renewcommand{\externaldocument}[2][]{%
  \IfFileExists{tmp/xref/#2.aux}{%
    \externaldocumentorig[#1]{tmp/xref/#2}[#2.pdf]%
  }{%
  \IfFileExists{#2.aux}{%
    \externaldocumentorig[#1]{#2}%
  }{%
    \PackageWarning{stack}{Missing #2.aux; external references with prefix #1 unavailable}%
  }%
  }%
}

% Package for hypertext links:
\usepackage{hyperref}

% For any local file, say "hello.tex" you want to link to please
% use \externaldocument[hello-]{hello}
\externaldocument[introduction-]{introduction}
\externaldocument[conventions-]{conventions}
\externaldocument[sets-]{sets}
\externaldocument[categories-]{categories}
\externaldocument[topology-]{topology}
\externaldocument[sheaves-]{sheaves}
\externaldocument[sites-]{sites}
\externaldocument[stacks-]{stacks}
\externaldocument[fields-]{fields}
\externaldocument[algebra-]{algebra}
\externaldocument[brauer-]{brauer}
\externaldocument[homology-]{homology}
\externaldocument[derived-]{derived}
\externaldocument[simplicial-]{simplicial}
\externaldocument[more-algebra-]{more-algebra}
\externaldocument[smoothing-]{smoothing}
\externaldocument[modules-]{modules}
\externaldocument[sites-modules-]{sites-modules}
\externaldocument[injectives-]{injectives}
\externaldocument[cohomology-]{cohomology}
\externaldocument[sites-cohomology-]{sites-cohomology}
\externaldocument[dga-]{dga}
\externaldocument[dpa-]{dpa}
\externaldocument[sdga-]{sdga}
\externaldocument[hypercovering-]{hypercovering}
\externaldocument[schemes-]{schemes}
\externaldocument[constructions-]{constructions}
\externaldocument[properties-]{properties}
\externaldocument[morphisms-]{morphisms}
\externaldocument[coherent-]{coherent}
\externaldocument[divisors-]{divisors}
\externaldocument[limits-]{limits}
\externaldocument[varieties-]{varieties}
\externaldocument[topologies-]{topologies}
\externaldocument[descent-]{descent}
\externaldocument[perfect-]{perfect}
\externaldocument[more-morphisms-]{more-morphisms}
\externaldocument[flat-]{flat}
\externaldocument[groupoids-]{groupoids}
\externaldocument[more-groupoids-]{more-groupoids}
\externaldocument[etale-]{etale}
\externaldocument[chow-]{chow}
\externaldocument[intersection-]{intersection}
\externaldocument[pic-]{pic}
\externaldocument[weil-]{weil}
\externaldocument[adequate-]{adequate}
\externaldocument[dualizing-]{dualizing}
\externaldocument[duality-]{duality}
\externaldocument[discriminant-]{discriminant}
\externaldocument[derham-]{derham}
\externaldocument[local-cohomology-]{local-cohomology}
\externaldocument[algebraization-]{algebraization}
\externaldocument[curves-]{curves}
\externaldocument[resolve-]{resolve}
\externaldocument[models-]{models}
\externaldocument[functors-]{functors}
\externaldocument[equiv-]{equiv}
\externaldocument[pione-]{pione}
\externaldocument[etale-cohomology-]{etale-cohomology}
\externaldocument[proetale-]{proetale}
\externaldocument[relative-cycles-]{relative-cycles}
\externaldocument[more-etale-]{more-etale}
\externaldocument[trace-]{trace}
\externaldocument[crystalline-]{crystalline}
\externaldocument[spaces-]{spaces}
\externaldocument[spaces-properties-]{spaces-properties}
\externaldocument[spaces-morphisms-]{spaces-morphisms}
\externaldocument[decent-spaces-]{decent-spaces}
\externaldocument[spaces-cohomology-]{spaces-cohomology}
\externaldocument[spaces-limits-]{spaces-limits}
\externaldocument[spaces-divisors-]{spaces-divisors}
\externaldocument[spaces-over-fields-]{spaces-over-fields}
\externaldocument[spaces-topologies-]{spaces-topologies}
\externaldocument[spaces-descent-]{spaces-descent}
\externaldocument[spaces-perfect-]{spaces-perfect}
\externaldocument[spaces-more-morphisms-]{spaces-more-morphisms}
\externaldocument[spaces-flat-]{spaces-flat}
\externaldocument[spaces-groupoids-]{spaces-groupoids}
\externaldocument[spaces-more-groupoids-]{spaces-more-groupoids}
\externaldocument[bootstrap-]{bootstrap}
\externaldocument[spaces-pushouts-]{spaces-pushouts}
\externaldocument[spaces-chow-]{spaces-chow}
\externaldocument[groupoids-quotients-]{groupoids-quotients}
\externaldocument[spaces-more-cohomology-]{spaces-more-cohomology}
\externaldocument[spaces-simplicial-]{spaces-simplicial}
\externaldocument[spaces-duality-]{spaces-duality}
\externaldocument[formal-spaces-]{formal-spaces}
\externaldocument[restricted-]{restricted}
\externaldocument[spaces-resolve-]{spaces-resolve}
\externaldocument[formal-defos-]{formal-defos}
\externaldocument[defos-]{defos}
\externaldocument[cotangent-]{cotangent}
\externaldocument[examples-defos-]{examples-defos}
\externaldocument[algebraic-]{algebraic}
\externaldocument[examples-stacks-]{examples-stacks}
\externaldocument[stacks-sheaves-]{stacks-sheaves}
\externaldocument[criteria-]{criteria}
\externaldocument[artin-]{artin}
\externaldocument[quot-]{quot}
\externaldocument[stacks-properties-]{stacks-properties}
\externaldocument[stacks-morphisms-]{stacks-morphisms}
\externaldocument[stacks-limits-]{stacks-limits}
\externaldocument[stacks-cohomology-]{stacks-cohomology}
\externaldocument[stacks-perfect-]{stacks-perfect}
\externaldocument[stacks-introduction-]{stacks-introduction}
\externaldocument[stacks-more-morphisms-]{stacks-more-morphisms}
\externaldocument[stacks-geometry-]{stacks-geometry}
\externaldocument[moduli-]{moduli}
\externaldocument[moduli-curves-]{moduli-curves}
\externaldocument[examples-]{examples}
\externaldocument[exercises-]{exercises}
\externaldocument[guide-]{guide}
\externaldocument[desirables-]{desirables}
\externaldocument[coding-]{coding}
\externaldocument[obsolete-]{obsolete}
\externaldocument[fdl-]{fdl}
\externaldocument[index-]{index}

% Heap Project documents
\externaldocument[k3-]{k3}
\externaldocument[torelli-]{torelli}
\externaldocument[linear-algebra-]{linear-algebra}
\externaldocument[groups-]{groups}
\externaldocument[complex-analysis-]{complex-analysis}
\externaldocument[functional-analysis-]{functional-analysis}
\externaldocument[categories-]{categories}
\externaldocument[commutative-algebra-]{commutative-algebra}
\externaldocument[ringed-spaces-]{ringed-spaces}
\externaldocument[homological-algebra-]{homological-algebra}
\externaldocument[lie-algebras-]{lie-algebras}
\externaldocument[representation-theory-]{representation-theory}
\externaldocument[smooth-manifolds-]{smooth-manifolds}
\externaldocument[differential-topology-]{differential-topology}
\externaldocument[algebraic-topology-]{algebraic-topology}
\externaldocument[differential-geometry-]{differential-geometry}
\externaldocument[algebraic-geometry-]{algebraic-geometry}
\externaldocument[symplectic-geometry-]{symplectic-geometry}
\externaldocument[complex-geometry-]{complex-geometry}
\externaldocument[hodge-theory-]{hodge-theory}
\externaldocument[%
variations-of-hodge-structures-]{%
variations-of-hodge-structures}
\externaldocument[%
homological-mirror-symmetry-]{%
homological-mirror-symmetry}
\maybeexternaldocument{deformations-}{deformations}
\externaldocument[vertex-algebras-]{vertex-algebras}
\externaldocument[springer-]{springer}
\externaldocument[infty-categories-]{infty-categories}
\externaldocument[seminars-differential-geometry-]{%
seminars-differential-geometry}
\externaldocument[seminars-complex-geometry-]{%
seminars-complex-geometry}
\externaldocument[seminars-algebraic-geometry-]{%
seminars-algebraic-geometry}
\externaldocument[seminars-analysis-pde-]{%
seminars-analysis-pde}
\externaldocument[surfaces-]{surfaces}
\externaldocument[birational-maps-commutative-algebra-]{birational-maps-commutative-algebra}
\externaldocument[cremona-transformations-]{cremona-transformations}
\externaldocument[derived-categories-of-sheaves-]{derived-categories-of-sheaves}
\externaldocument[geometric-invariant-theory-]{geometric-invariant-theory}
\externaldocument[geometric-stability-conditions-and-group-actions-]{geometric-stability-conditions-and-group-actions}
\externaldocument[moduli-spaces-of-sheaves-]{moduli-spaces-of-sheaves}
\externaldocument[bridgeland-stability-general-theory-]{bridgeland-stability-general-theory}
\externaldocument[hilbert-schemes-]{hilbert-schemes}
\externaldocument[noncommutative-abelian-surfaces-and-kummer-type-hyperkahler-varieties-]{noncommutative-abelian-surfaces-and-kummer-type-hyperkahler-varieties}
\externaldocument[mumford-tate-groups-in-hodge-theory-]{mumford-tate-groups-in-hodge-theory}
\externaldocument[hodge-birational-atoms-2026-]{hodge-birational-atoms-2026}
\externaldocument[xxii-egd-2026-]{%
xxii-egd-2026}
\externaldocument[pde-]{pde}
\externaldocument[probability-]{probability}

\externaldocument[linguistics-]{linguistics}

\externaldocument[%
16th-alga-meeting-2026-]{%
16th-alga-meeting-2026}

\externaldocument[%
iv-jornadas-lefschetz-2026-]{%
iv-jornadas-lefschetz-2026}

\externaldocument[reading-the-masters-]{%
reading-the-masters}

\externaldocument[real-analysis-]{%
real-analysis}

\externaldocument[optimal-transport-]{%
optimal-transport}

\externaldocument[history-]{history}

\externaldocument[seiberg-witten-]{seiberg-witten}

\externaldocument[history-philosophy-biology-]{history-philosophy-biology}

\externaldocument[classical-mechanics-]{classical-mechanics}
\externaldocument[electromagnetism-]{electromagnetism}
\externaldocument[waves-and-optics-]{waves-and-optics}
\externaldocument[wif-]{waves-in-fluids-iv}
\externaldocument[thermal-physics-]{thermal-physics}
\externaldocument[quantum-mechanics-]{quantum-mechanics}
\externaldocument[special-relativity-]{special-relativity}
\externaldocument[physics-]{physics}

\externaldocument[chemistry-]{chemistry}

\externaldocument[molecular-cell-biology-]{molecular-cell-biology}

\externaldocument[microbiology-]{microbiology}

\externaldocument[protistology-]{protistology}

\externaldocument[mycology-]{mycology}

\externaldocument[botany-]{botany}

\externaldocument[zoology-]{zoology}

\externaldocument[animal-histology-]{animal-histology}

\externaldocument[animal-physiology-]{animal-physiology}

\externaldocument[systematics-biogeography-]{systematics-biogeography}

\externaldocument[ecology-natural-resources-]{ecology-natural-resources}

\externaldocument[earth-sciences-]{earth-sciences}

\externaldocument[biostatistics-]{biostatistics}

\externaldocument[mathematical-systems-biology-]{mathematical-systems-biology}

\externaldocument[cimpa-el-salvador-2024-]{cimpa-el-salvador-2024}

\externaldocument[anthropology-]{anthropology}

\externaldocument[yang-mills-theory-]{yang-mills-theory}

\externaldocument[seminars-algebra-]{seminars-algebra}

\externaldocument[seminars-physics-]{seminars-physics}

\externaldocument[seminars-computing-]{seminars-computing}

\externaldocument[oktobermat-2026-]{oktobermat-2026}

\externaldocument[saet-2026-]{saet-2026}

\externaldocument[seminars-symplectic-geometry-]{seminars-symplectic-geometry}

\externaldocument[seminars-dynamics-]{seminars-dynamics}

\externaldocument[oktobermat-2024-]{oktobermat-2024}

\externaldocument[quantum-groups-]{quantum-groups}

\externaldocument[legal-2027-]{legal-2027}

\externaldocument[derived-categories-]{derived-categories}

% Theorem environments.
%
\theoremstyle{plain}
\newtheorem{theorem}[subsection]{Theorem}
\newtheorem{proposition}[subsection]{Proposition}
\newtheorem{lemma}[subsection]{Lemma}
\newtheorem{corollary}[subsection]{Corollary}

\theoremstyle{definition}
\newtheorem{definition}[subsection]{Definition}
\newtheorem{example}[subsection]{Example}
\newtheorem{exercise}[subsection]{Exercise}
\newtheorem{situation}[subsection]{Situation}

\theoremstyle{remark}
\newtheorem{remark}[subsection]{Remark}
\newtheorem{remarks}[subsection]{Remarks}

\numberwithin{equation}{subsection}

% Macros
%
\def\lim{\mathop{\mathrm{lim}}\nolimits}
\def\colim{\mathop{\mathrm{colim}}\nolimits}
\def\Spec{\mathop{\mathrm{Spec}}}
\def\Hom{\mathop{\mathrm{Hom}}\nolimits}
\def\Ext{\mathop{\mathrm{Ext}}\nolimits}
\def\SheafHom{\mathop{\mathcal{H}\!\mathit{om}}\nolimits}
\def\SheafExt{\mathop{\mathcal{E}\!\mathit{xt}}\nolimits}
\def\Sch{\mathit{Sch}}
\def\Mor{\mathop{\mathrm{Mor}}\nolimits}
\def\Ob{\mathop{\mathrm{Ob}}\nolimits}
\def\Sh{\mathop{\mathit{Sh}}\nolimits}
\def\NL{\mathop{N\!L}\nolimits}
\def\CH{\mathop{\mathrm{CH}}\nolimits}
\def\proetale{{pro\text{-}\acute{e}tale}}
\def\etale{{\acute{e}tale}}
\def\QCoh{\mathit{QCoh}}
\def\Ker{\mathop{\mathrm{Ker}}}
\def\Im{\mathop{\mathrm{Im}}}
\def\Coker{\mathop{\mathrm{Coker}}}
\def\Coim{\mathop{\mathrm{Coim}}}

% Boxtimes
%
\DeclareMathSymbol{\boxtimes}{\mathbin}{AMSa}{"02}

%
% Macros for moduli stacks/spaces
%
\def\QCohstack{\mathcal{QC}\!\mathit{oh}}
\def\Cohstack{\mathcal{C}\!\mathit{oh}}
\def\Spacesstack{\mathcal{S}\!\mathit{paces}}
\def\Quotfunctor{\mathrm{Quot}}
\def\Hilbfunctor{\mathrm{Hilb}}
\def\Curvesstack{\mathcal{C}\!\mathit{urves}}
\def\Polarizedstack{\mathcal{P}\!\mathit{olarized}}
\def\Complexesstack{\mathcal{C}\!\mathit{omplexes}}
% \Pic is the operator that assigns to X its picard group, usage \Pic(X)
% \Picardstack_{X/B} denotes the Picard stack of X over B
% \Picardfunctor_{X/B} denotes the Picard functor of X over B
\def\Pic{\mathop{\mathrm{Pic}}\nolimits}
\def\Picardstack{\mathcal{P}\!\mathit{ic}}
\def\Picardfunctor{\mathrm{Pic}}
\def\Deformationcategory{\mathcal{D}\!\mathit{ef}}


\begin{document}

\title{Derived categories}
\maketitle

\phantomsection
\label{section-phantom}

\tableofcontents

\section{Overview}
\label{section-overview}

\noindent
These notes start from the lecture of
23 November 2026 in Jethro van Ekeren's
course on homological algebra (IMPA).
They are meant as a skeleton: the steps
of the construction and why each one is
taken, without the full proofs.

\medskip\noindent
{\bf Sources.} My own notes of the
lecture; two photographs of the board;
spoken remarks, paraphrased; Sections
9.1 and 9.2 of Jethro's
\href{https://w3.impa.br/~jethro/2026-2/%
ah-current.pdf}{course notes}, which say
they are based on
\cite[Chapter 2]{huybrechts-fm}.
Statements checked against
\cite{huybrechts-fm} carry the exact
number. Paragraphs marked
{\it editorial} were added afterwards
and were not said in the lecture;
{\it from the lecture} marks remarks that
were said and not written on the board.
A minicourse on the same subject is in
Section
\ref{derived-categories-of-sheaves-section-motivation}
and the sections around it.

\medskip\noindent
{\bf Big picture.} Homological algebra
replaces an object by a complex and
only cares about the complex up to
something weaker than isomorphism.
The derived category is the place where
this ``up to'' becomes literal
isomorphism of objects. It is built in
steps:
\begin{enumerate}
\item Start with an abelian category
$\mathcal{A}$.
\item Pass to the category of complexes
$\operatorname{Ch}(\mathcal{A})$. It is
abelian.
\item Pass to the homotopy category
$K(\mathcal{A})$: same objects, maps up
to homotopy. It is additive but {\bf not}
abelian in general.
\item Since there are no exact sequences
in $K(\mathcal{A})$, replace them by
{\it distinguished triangles}, which come
from cones.
\item Prove that $K(\mathcal{A})$ with
these triangles is a
{\it triangulated category}.
\item Only now write down the definition
of triangulated category: its axioms are
the properties just observed.
\item (Next.) Invert the
quasi-isomorphisms of $K(\mathcal{A})$
to get the derived category
$D(\mathcal{A})$.
\end{enumerate}

\section{Complexes and the homotopy category}
\label{section-complexes-homotopy-category}

\noindent
Let $\mathcal{A}$ be an abelian category.
Jethro always works with cohomological
complexes
$$
\xymatrix{
\cdots\ar[r]
&A^{-1}\ar[r]^{d}
&A^{0}\ar[r]^{d}
&A^{1}\ar[r]
&\cdots
}
$$
with $d\circ d=0$. The category of
complexes is written
$\operatorname{Ch}(\mathcal{A})$, or
$\operatorname{Com}(\mathcal{A})$, or
$\operatorname{Kom}(\mathcal{A})$,
depending on the author
(the last one is the notation of
\cite{huybrechts-fm}). It is again
abelian, see Lemma
\ref{homological-algebra-lemma-cat-cochain-abelian}.

\medskip\noindent
{\bf From the board.} A short exact
sequence
$$
\xymatrix{
0\ar[r]&A\ar[r]&B\ar[r]&C\ar[r]&0
}
$$
in $\operatorname{Ch}(\mathcal{A})$ gives
a long exact sequence in $\mathcal{A}$
$$
\xymatrix@C=1.1em{
\cdots\ar[r]
&H^n(A)\ar[r]
&H^n(B)\ar[r]
&H^n(C)\ar[r]
&H^{n+1}(A)\ar[r]
&\cdots
}
$$
Two chain maps $f,g:A\to B$ are
{\it homotopic}, written $f\sim g$, if
there are maps $h:A^n\to B^{n-1}$ with
$$
f-g=dh+hd.
$$
``Marvellous fact'': if $f\sim g$ then
$$
H^n(f)=H^n(g)
$$
for all $n$ (Lemma
\ref{homological-algebra-lemma-map-cohomology-homotopy-cochain}).
If $f:A\to B$ and $g:B\to A$ satisfy
$f\circ g\sim\operatorname{id}$ and
$g\circ f\sim\operatorname{id}$, then
$H^n(A)\cong H^n(B)$ for all $n$; such
$f$ and $g$ are called homotopy
equivalences. A chain map $f:A\to B$ for
which every
$$
H^n(f):H^n(A)\to H^n(B)
$$
is an isomorphism is called a
quasi-isomorphism (``qis'' for short).
The board had a picture of a small
blob of homotopy equivalences inside a
big blob of quasi-isomorphisms: every
homotopy equivalence is a
quasi-isomorphism, but not conversely.

\begin{definition}
\label{definition-homotopy-category}
The {\it homotopy category}
$K(\mathcal{A})$ has the same objects as
$\operatorname{Ch}(\mathcal{A})$, and
for two complexes $A$ and $B$
$$
\operatorname{Hom}_{K(\mathcal{A})}(A,B)
=
\operatorname{Hom}_{
\operatorname{Ch}(\mathcal{A})}(A,B)
/\sim.
$$
\end{definition}

\noindent
Some authors write
$\mathcal{H}(\mathcal{A})$. The
category $K(\mathcal{A})$ is still
additive, but it is not abelian in
general: there are morphisms with no
kernel (an example is in Remark
\ref{remark-homotopy-category-not-abelian}).
Complexes related by a homotopy
equivalence become isomorphic in
$K(\mathcal{A})$.

\medskip\noindent
{\it From the lecture.} The quotient makes sense
because being homotopic to zero is
``kind of an ideal'' in the set of all
morphisms: composing homotopic maps with
any map gives homotopic maps (Lemma
\ref{homological-algebra-lemma-compose-homotopy}).
The definition is simple but tricky to
get your head around: there are fewer
morphisms, since we took a quotient, and
yet more isomorphisms, since it became
easier for two objects to be isomorphic.

\section{Shift, triangles and cones}
\label{section-shift-triangles-cones}

\noindent
Both Jethro and Lyalya Guseva (in the
minicourse) introduced the shift at this
point.

\begin{definition}
\label{definition-shift}
The {\it shift} (or {\it suspension})
of a complex $A$ is the complex $A[1]$
with
$$
A[1]^n=A^{n+1},
\qquad
d_{A[1]}=-d_A.
$$
For a chain map $f:A\to B$, the map
$f[1]:A[1]\to B[1]$ is
$f[1]^n=f^{n+1}$.
\end{definition}

\noindent
So everything moves one place to the
left. The minus sign has to appear
somewhere, whatever the convention; where
it appears is a choice. We will see it
again in Proposition
\ref{proposition-huybrechts-2-16} and in
axiom (TR2).

\medskip\noindent
{\bf Why triangles.} We cannot take
exact sequences in $K(\mathcal{A})$,
because it is not abelian. What replaces
them is the following.

\begin{definition}
\label{definition-triangle}
Let $\mathcal{C}$ be an additive category
with an equivalence
$\Sigma:\mathcal{C}\to\mathcal{C}$,
written $\Sigma X=X[1]$.
A {\it triangle} in $\mathcal{C}$ is a
diagram
$$
\xymatrix{
A\ar[r]^{f}
&B\ar[r]^{g}
&C\ar[r]^{h}
&A[1]
}
$$
\end{definition}

\noindent
{\bf The cone gives triangles.} We
already used the cone in the course,
see Section
\ref{homological-algebra-section-cone-construction}.
Let $f:A\to B$ be a morphism in
$\operatorname{Ch}(\mathcal{A})$. Its
cone $C(f)$ is
$$
C(f)^n=A^{n+1}\oplus B^n,
$$
$$
d_{C(f)}(a,b)=(-d_A(a),f(a)+d_B(b)).
$$
Writing elements as columns, this is
the matrix
$$
d_{C(f)}=
\begin{pmatrix}
-d_A&0\\
f&d_B
\end{pmatrix}
$$
of \cite[Definition 2.15]{huybrechts-fm}.
(In my notes I wrote the transpose,
which is the same map acting on rows.)
There are two natural maps: the
inclusion $\tau:B\to C(f)$ and the
projection $\pi:C(f)\to A[1]$. By
Proposition
\ref{homological-algebra-proposition-ses-complexes-cone}
they form a short exact sequence of
complexes
\begin{equation}
\label{equation-ses-cone}
\xymatrix{
0\ar[r]
&B\ar[r]^-{\tau}
&C(f)\ar[r]^-{\pi}
&A[1]\ar[r]
&0
}
\end{equation}
and so a long exact sequence
$$
\xymatrix@C=0.9em{
\cdots\ar[r]
&H^i(B)\ar[r]
&H^i(C(f))\ar[r]
&H^{i+1}(A)\ar[r]
&H^{i+1}(B)\ar[r]
&\cdots
}
$$

\begin{definition}
\label{definition-cone-triangle}
The {\it triangle of the cone} of
$f:A\to B$ is
\begin{equation}
\label{equation-cone-triangle}
\xymatrix{
A\ar[r]^-{f}
&B\ar[r]^-{\tau}
&C(f)\ar[r]^-{\pi}
&A[1]
}
\end{equation}
A triangle in $K(\mathcal{A})$ is
{\it distinguished} if it is isomorphic
to the triangle of a cone.
\end{definition}

\noindent
{\bf Idea.} A distinguished triangle is
what is left of an exact sequence. It is
meant to model ``giving rise to a long
exact sequence''. Instead of ending in
zero, it returns to $A$, one step higher.

\medskip\noindent
{\it From the lecture.} In $K(\mathcal{A})$ it no
longer makes sense to say that $B$
injects into the middle term and the
middle term surjects onto the last one:
``where we are on the triangle loses
meaning''. What survives is the long
exact sequence, and a long exact
sequence has a symmetry: it can be
rotated. That is why one expects to be
able to rotate distinguished triangles.

\begin{remark}
\label{remark-triangle-is-helix}
(Jethro's notes, Remark 9.2.1.) A
triangle is really a helix: the triangle
\eqref{equation-cone-triangle} is no more
fundamental than its rotation
$$
\xymatrix{
B\ar[r]^-{\tau}
&C(f)\ar[r]^-{\pi}
&A[1]\ar[r]^-{-f[1]}
&B[1]
}
$$
The triangle \eqref{equation-cone-triangle}
is distinguished because
\eqref{equation-ses-cone} is exact. But
$A\to B\to C(f)$ is not exact, and the
composition $\tau\circ f$ is not even
zero in
$\operatorname{Ch}(\mathcal{A})$. It is
zero in $K(\mathcal{A})$ (Jethro's
Exercise 9.2.2).
\end{remark}

\medskip\noindent
{\it Editorial: solution of the
exercise.} Take $h:A^n\to C(f)^{n-1}$,
$h(a)=(a,0)$; note that
$C(f)^{n-1}=A^n\oplus B^{n-1}$. Then
$$
\begin{aligned}
dh(a)+hd(a)
&=(-d_A(a),f(a))+(d_A(a),0)\\
&=(0,f(a)),
\end{aligned}
$$
which is $\tau(f(a))$. So
$\tau\circ f\sim 0$.

\section{The homotopy category is triangulated}
\label{section-homotopy-category-triangulated}

\noindent
We want to define a triangulated
category: an additive category
$\mathcal{C}$ with a shift functor
$\Sigma$ and a collection of
distinguished triangles, subject to
axioms. Before writing the axioms we
check, in $K(\mathcal{A})$, the facts
that the axioms will ask for.

\begin{lemma}
\label{lemma-cone-identity-zero}
For every complex $A$, the cone of the
identity is isomorphic to zero in
$K(\mathcal{A})$:
$$
C(\operatorname{id}_A)\cong 0.
$$
\end{lemma}

\begin{proof}
We need a homotopy from the identity of
$C(\operatorname{id}_A)$ to zero. Here
$$
C(\operatorname{id}_A)^n
=A^{n+1}\oplus A^n,
$$
$$
d(a,b)=(-d_A(a),a+d_A(b)).
$$
Take $h(a,b)=(b,0)$, which lands in
$A^n\oplus A^{n-1}$. Then
$$
\begin{aligned}
dh(a,b)&=(-d_A(b),b),\\
hd(a,b)&=(a+d_A(b),0),
\end{aligned}
$$
and the sum is $(a,b)$. (Jethro's notes
leave this check as an exercise; the
computation here is editorial.)
\end{proof}

\noindent
So
$$
\xymatrix{
A\ar[r]^-{\operatorname{id}}
&A\ar[r]
&0\ar[r]
&A[1]
}
$$
is a distinguished triangle. This will
be axiom (TR1)(i).

\medskip\noindent
{\bf Commuting squares.} A commuting
square
$$
\xymatrix{
A\ar[r]^{f}\ar[d]_{\alpha}
&B\ar[d]^{\beta}\\
A'\ar[r]^{f'}
&B'
}
$$
gives a map $C(f)\to C(f')$, namely
$\alpha\oplus\beta$ in each degree, and
with it a morphism between the two cone
triangles. Whether the squares of this
morphism commute on the nose or only up
to homotopy, and what happens if the
first square commutes only up to
homotopy, is Jethro's Exercise 9.2.3.
This will be axiom (TR3).

\medskip\noindent
{\bf Rotation.} This is where the real
work is. We want: if
$$
\xymatrix{
A\ar[r]^{f}
&B\ar[r]^{g}
&C\ar[r]^{h}
&A[1]
}
$$
is distinguished, then so is
$$
\xymatrix{
B\ar[r]^{g}
&C\ar[r]^{h}
&A[1]\ar[r]^-{-f[1]}
&B[1]
}
$$
Since distinguished means ``isomorphic
to a cone triangle'', we must compare
the rotated triangle with the cone
triangle of $\tau:B\to C(f)$. This is
the content of:

\begin{proposition}
\label{proposition-huybrechts-2-16}
\cite[Proposition 2.16]{huybrechts-fm}
Let $f:A\to B$ be a morphism of
complexes, with $\tau:B\to C(f)$ and
$\pi:C(f)\to A[1]$ as above. Then there
is a morphism of complexes
$$
g:A[1]\to C(\tau)
$$
which is an isomorphism in
$K(\mathcal{A})$ and such that the
diagram
$$
\xymatrix{
B\ar[r]^-{\tau}\ar@{=}[d]
&C(f)\ar[r]^-{\pi}\ar@{=}[d]
&A[1]\ar[r]^-{-f}\ar[d]^{g}
&B[1]\ar@{=}[d]\\
B\ar[r]^-{\tau}
&C(f)\ar[r]
&C(\tau)\ar[r]
&B[1]
}
$$
commutes in $K(\mathcal{A})$. The bottom
row is the cone triangle of $\tau$.
\end{proposition}

\noindent
{\bf The map.} We have
$$
C(\tau)^i
=B^{i+1}\oplus C(f)^i
=B^{i+1}\oplus A^{i+1}\oplus B^i
$$
and $A[1]^i=A^{i+1}$. The map is
$$
g=(-f^{i+1},\operatorname{id},0).
$$
Its inverse in $K(\mathcal{A})$ is the
projection to the middle factor. The
right square commutes already in
$\operatorname{Ch}(\mathcal{A})$; the
middle square commutes only up to
homotopy \cite[proof of Proposition
2.16]{huybrechts-fm}. The sign in $g$ is
the reason for the sign $-f[1]$ in the
rotated triangle.

\begin{theorem}
\label{theorem-homotopy-category-triangulated}
(Jethro's notes, Theorem 9.2.4;
\cite[Proposition 2.24]{huybrechts-fm}.)
The category $K(\mathcal{A})$, with the
shift of complexes and with
distinguished triangles those isomorphic
to the triangle of a cone, is a
triangulated category in the sense of
Definition
\ref{definition-triangulated-category}.
\end{theorem}

\noindent
{\bf Status.} What was checked above is
(TR1), (TR3) and one direction of (TR2).
Jethro said that he feels a more general
version of the argument of Proposition
\ref{proposition-huybrechts-2-16} gives
both (TR2) and (TR3). The octahedral
axiom (TR4) was not verified (``it can
be done''); Jethro's notes say the
theorem is ``essentially verified'' if
one ignores it, and \cite{huybrechts-fm}
refers to the literature for the proof.

\section{Triangulated categories}
\label{section-triangulated-categories}

\noindent
Only now comes the definition. Each axiom
is one of the facts of Section
\ref{section-homotopy-category-triangulated}.

\begin{definition}
\label{definition-triangulated-category}
A {\it triangulated category} is an
additive category $\mathcal{C}$ with an
equivalence
$\Sigma:\mathcal{C}\to\mathcal{C}$,
written $\Sigma X=X[1]$, and a collection
of triangles called {\it distinguished},
such that:
\begin{enumerate}
\item[(TR1)] (i) For every $A$ the
triangle
$$
\xymatrix{
A\ar[r]^-{\operatorname{id}}
&A\ar[r]
&0\ar[r]
&A[1]
}
$$
is distinguished. (ii) A triangle
isomorphic to a distinguished triangle is
distinguished. (iii) Every morphism
$f:A\to B$ is the first map of some
distinguished triangle
$$
\xymatrix{
A\ar[r]^{f}
&B\ar[r]^{g}
&C\ar[r]^{h}
&A[1]
}
$$
\item[(TR2)] The triangle
$$
\xymatrix{
A\ar[r]^{f}
&B\ar[r]^{g}
&C\ar[r]^{h}
&A[1]
}
$$
is distinguished if and only if
$$
\xymatrix{
B\ar[r]^{g}
&C\ar[r]^{h}
&A[1]\ar[r]^-{-f[1]}
&B[1]
}
$$
is distinguished.
\item[(TR3)] Given two distinguished
triangles and maps $\alpha$, $\beta$
making the left square commute,
$$
\xymatrix{
A\ar[r]^{f}\ar[d]_{\alpha}
&B\ar[r]\ar[d]^{\beta}
&C\ar[r]\ar@{-->}[d]^{\gamma}
&A[1]\ar[d]^{\alpha[1]}\\
A'\ar[r]^{f'}
&B'\ar[r]
&C'\ar[r]
&A'[1]
}
$$
there is a map $\gamma$ making the whole
diagram commute.
\item[(TR4)] (Octahedral axiom.) Given
two composable morphisms
$$
\xymatrix{
A\ar[r]^{f}
&B\ar[r]^{g}
&C
}
$$
and distinguished triangles on $f$, on
$g$ and on $g\circ f$, with third objects
$C'$, $A'$ and $B'$ respectively, there
is a distinguished triangle
$$
\xymatrix{
C'\ar[r]
&B'\ar[r]
&A'\ar[r]
&C'[1]
}
$$
compatible with the maps of the three
given triangles.
\end{enumerate}
\end{definition}

\noindent
This is Definition 9.2.2 of Jethro's
notes; compare
\cite[Definition 1.32]{huybrechts-fm}.
The compatibility in (TR4) is not
spelled out here: Jethro's notes leave
that sentence unfinished, and
\cite{huybrechts-fm} does not state
(TR4) precisely either, saying that it
is never used explicitly in the book.

\begin{remark}
\label{remark-octahedral-intuition}
(Idea behind the octahedral axiom;
Jethro's notes, Remark 9.2.5; the same
explanation follows
\cite[Definition 1.32]{huybrechts-fm}.)
Take abelian groups
$A\subset B\subset C$. Then
$$
C/B\cong (C/A)/(B/A),
$$
that is, there is a short exact sequence
$$
\xymatrix{
0\ar[r]
&B/A\ar[r]
&C/A\ar[r]
&C/B\ar[r]
&0
}
$$
Its three terms come from three short
exact sequences of the original setup:
$$
\xymatrix@R=0.6em{
0\ar[r]&A\ar[r]&B\ar[r]&B/A\ar[r]&0\\
0\ar[r]&B\ar[r]&C\ar[r]&C/B\ar[r]&0\\
0\ar[r]&A\ar[r]&C\ar[r]&C/A\ar[r]&0
}
$$
Now replace ``short exact sequence'' by
``distinguished triangle''. The
inclusions become $f:A\to B$ and
$g:B\to C$; the quotients $B/A$, $C/B$
and $C/A$ become the third objects $C'$,
$A'$ and $B'$ of triangles on $f$, on
$g$ and on $g\circ f$. The axiom says
that the fourth sequence also exists:
$C'\to B'\to A'\to C'[1]$ is
distinguished, in the same way that
$B/A\to C/A\to C/B$ is exact.
\end{remark}

\noindent
\textbf{Editorial note.} The end of the
remark is the dictionary in its last
paragraph: a quotient by a subobject
becomes the third vertex of a triangle.
The octahedral axiom is then the third
isomorphism theorem for cones: ``the cone
of $g\circ f$ modulo the cone of $f$ is
the cone of $g$''. For $K(\mathcal{A})$
one can see the maps: Jethro's notes
point out that for composable $f$ and
$g$ there are canonical morphisms
$$
\xymatrix{
C(f)\ar[r]
&C(g\circ f)\ar[r]
&C(g)
}
$$
given in degree $i$ by
$$
(a,b)\mapsto(a,g(b)),
\qquad
(a,c)\mapsto(f(a),c).
$$
The explicit formulas are editorial;
that both are chain maps is a direct
check with the differential of the cone.

\begin{quote}
\textbf{Question.} ``Is it true
that the octahedral axiom was recently
proved to actually be necessary (like
Euclid's 5th postulate is indeed
necessary)? Do we get non-canonical
derived categories if we deny the
octahedral axiom?''
\end{quote}

\noindent
\textbf{Answer.} Yes to the
first question, according to a 2026
preprint. In \cite{cilsz-tr4} the
authors construct classes of triangles
satisfying (TR1), (TR2) and (TR3) but
not (TR4); their abstract says this
``confirms a long-standing conjecture on
triangulated categories''. So (TR4) does
not follow from the other axioms. Only
the abstract was read; it is a preprint
(arXiv:2608.09777), not yet a refereed
paper.

\medskip\noindent
To the second question: not in the way
the analogy suggests. By its abstract,
\cite{cilsz-tr4} perturbs the class of
distinguished triangles on a category
that is already triangulated in the
usual way. For $K(\mathcal{A})$ and
$D(\mathcal{A})$ with the cone triangles
(TR4) is a theorem, not a choice
\cite[Proposition 2.24]{huybrechts-fm},
so an abelian category has no
``non-octahedral derived category''.
Whether the exotic structures have a
meaning of their own, as non-Euclidean
geometries do, I do not know.

\begin{remark}
\label{remark-homotopy-category-not-abelian}
(Editorial, from Jethro's notes, Section
9.4; not in this lecture.) In a
triangulated category every
monomorphism splits. This gives a
morphism with no kernel in
$K(\mathcal{A})$ for $\mathcal{A}$ the
category of abelian groups: reduction
modulo $p$
$$
\mathbb{Z}/p^2\to\mathbb{Z}/p,
$$
both placed in degree $0$. So
triangulated categories are ``bad''
(not abelian) for the same reason that
they are ``good'' (all monomorphisms
split).
\end{remark}

\section{Next: the derived category}
\label{section-next-derived-category}

\noindent
(Said in the last minutes of the
lecture as a preview; written here with
the help of Jethro's notes, Section 9.3.
See also Definition
\ref{derived-categories-of-sheaves-definition-derived-category}.)
Passing from
$\operatorname{Ch}(\mathcal{A})$ to
$K(\mathcal{A})$ turned homotopy
equivalences into isomorphisms. The
derived category $D(\mathcal{A})$ does
the same for all quasi-isomorphisms. The
reason to want this: two projective
resolutions $P_1$ and $P_2$ of an object
$X$ are homotopy equivalent, so
$P_1\cong P_2$ already in
$K(\mathcal{A})$; but the maps
$P_1\to X$ and $P_2\to X$ are only
quasi-isomorphisms. In $D(\mathcal{A})$
one has
$$
P_1\cong X\cong P_2.
$$
A morphism from $A$ to $B$ in
$D(\mathcal{A})$ is a class of ``roofs''
$$
\xymatrix{
&C\ar[dl]_{s}\ar[dr]^{f}&\\
A&&B
}
$$
with $s$ a quasi-isomorphism, thought of
as the fraction of $f$ by $s$. Composing
two roofs needs
\cite[Proposition 2.17]{huybrechts-fm},
whose proof uses the cone and
Proposition
\ref{proposition-huybrechts-2-16}. That
is why the triangulated structure of
$K(\mathcal{A})$ comes first.

\medskip\noindent
{\it From the lecture.} Every chain map $f$ gives
a roof, with the identity as its
quasi-isomorphism leg. Every
quasi-isomorphism $q:B\to A$ gives a roof
from $A$ to $B$, with $q$ on one leg and
the identity on the other; this is the
inverse of $q$. Checking that it really
is an inverse, and that composition is
associative, works only up to homotopy,
which is why one passes through
$K(\mathcal{A})$. One can reach
$D(\mathcal{A})$ without the homotopy
category, but that takes heavy category
theory, and to compute anything one has
to bring the homotopy category back.

\bibliography{my}
\bibliographystyle{amsalpha}

\end{document}
